\documentclass[12pt]{article}
\usepackage[T1]{fontenc}
\usepackage[utf8]{inputenc}
\usepackage{lmodern}
\usepackage{textcomp}
\usepackage{enumitem}
\usepackage{geometry}
\usepackage{graphicx}
\usepackage{amsmath, amssymb}
\geometry{margin=1in}
\begin{document}
\begin{center}
Biology - Undergraduate Level Assessment 10 \
Total Marks: 40 \
Answer all questions.
\end{center}

\section*{Multiple Choice (10 marks)}
\begin{enumerate}[label=\arabic*.]
\item Which of the following is NOT a source of atmospheric carbon?
\begin{enumerate}[label=(\alph*)]
    \item Respiration
    \item Photosynthesis
    \item Bacterial decomposition
    \item Combustion of fossil fuels
\end{enumerate}
\item Which of the following depicts the correct sequence of membranes of the chloroplast, beginning with the innermost membrane and ending with the outermost membrane?
\begin{enumerate}[label=(\alph*)]
    \item Outer membrane, inner membrane, thylakoid membrane
    \item Thylakoid membrane, inner membrane, outer membrane
    \item Inner membrane, thylakoid membrane, outer membrane
    \item Thylakoid membrane, outer membrane, inner membrane
    \item Strama, outer membrane, inner membrane
\end{enumerate}
\item Which statement about variation is true?
\begin{enumerate}[label=(\alph*)]
    \item All phenotypic variation is the result of genotypic variation.
    \item All genetic variation produces phenotypic variation.
    \item All nucleotide variability results in neutral variation.
    \item All new alleles are the result of nucleotide variability.
\end{enumerate}
\item Most animal cells, regardless of species, are relatively small and about the same size. Relative to larger cells, why is this?
\begin{enumerate}[label=(\alph*)]
    \item Smaller cells are more resilient to external damage.
    \item Smaller cells have a larger number of organelles.
    \item Smaller cells require less energy to function.
    \item Smaller cells can multiply faster.
    \item Smaller cells fit together more tightly.
\end{enumerate}
\item Short sequence by promoter that assists transcription by interacting with regulatory proteins.
\begin{enumerate}[label=(\alph*)]
    \item Promoter
    \item Inducer
    \item Repressor
    \item Operator
    \item Inhibitor
\end{enumerate}
\end{enumerate}

\section*{True / False (10 marks)}
\textit{Answer True or False}
\begin{enumerate}[label=\arabic*.]
\setcounter{enumi}{5}
\item True or False: Calcitonin acts antagonistically to parathormone by lowering blood calcium levels.
\item True or False: The probability of rolling a 5 on two successive rolls of a balanced die is 1/24.
\item True or False: Flowers with strong sweet fragrances are primarily adapted to attract hummingbirds rather than bees.
\item True or False: Meiosis produces two haploid gametes after the completion of meiosis I.
\item True or False: Operant conditioning can be used to assess a dog's sensitivity to sounds of different frequencies through learned responses.
\end{enumerate}

\section*{Long Form Response (20 marks)}
\textit{Answer each question with a clear and complete explanation.}
\begin{enumerate}[label=\arabic*.]
\setcounter{enumi}{10}
\item Discuss the implications of removing a bacterial cell wall through lysozyme treatment. Can the bacteria survive, and how do bacterial and plant cell walls differ biochemically?
\item Discuss the role of plasma components, particularly inorganic salts, in influencing blood pressure. Additionally, analyze other solutes present in plasma and their potential effects.
\end{enumerate}

\vfill
\noindent\textit{End of Paper}
\end{document}